\documentclass[11pt]{article}

\usepackage[margin=1in]{geometry}
\usepackage{url}
\usepackage{xurl}
\usepackage{booktabs}
\usepackage{enumitem}
\usepackage{hyperref}
\usepackage{array}
\usepackage{tabularx}
\usepackage[numbers,sort&compress]{natbib}
\hypersetup{hidelinks,breaklinks=true}
\setlength{\emergencystretch}{2em}

\newcommand{\diagrambox}[2]{%
    \fbox{\parbox[c][4.6em][c]{#1}{\centering\footnotesize #2}}%
}

\title{
SkillOps: A Practical Framework for Designing, Testing, and Operating
Modular Skills in Personal AI Agents
}

\author{
Song Luo\\
\texttt{luosongred@gmail.com}
}

\date{April 2026}

\begin{document}

\maketitle

\begin{abstract}
Personal AI agents increasingly rely on reusable capability modules, often
called skills, yet their design and operation remain largely informal. This
paper introduces SkillOps, a practical framework for designing, testing, and
operating modular skills in personal AI agents. The framework is derived from
open-source artifacts developed by the author and contributes a skill component
taxonomy, a lifecycle model, a failure-mode taxonomy, a design-gap matrix, and
an artifact-based evaluation design. The paper reports three evidence layers:
descriptive benchmark summaries from manually constructed artifact, trigger,
and risk cases; deterministic repository-level consistency checks; and a local
security-guard pilot over risk and benign-control cases. Model-backed
experiments are specified as executable follow-up protocols, but no
model-backed live metrics are reported without the corresponding credentials
and run logs. The contribution is therefore a bounded engineering framework
for making agent skills more explicit, testable, and maintainable.
\end{abstract}

\section{Introduction}

Large language model agents are increasingly used as interactive systems that
can read instructions, call tools, edit files, use external APIs, and maintain
task-specific context. In such systems, user-visible behavior is shaped not
only by the base model, but also by auxiliary capability modules. These modules
may be called tools, plugins, agents, memories, workflows, or skills. This
paper focuses on the skill as a practical unit of agent capability in personal
agent environments.

In many personal agent settings, skills are written and maintained by
individual users or small teams rather than by centralized platform engineers.
A skill may contain a trigger description, operational instructions, examples,
constraints, memory rules, safety checks, and tests. When these elements are
implicit or poorly specified, the agent may activate the wrong skill, inject
excessive context, ignore constraints, retain obsolete information, or perform
unsafe operations. These are not only model-quality issues; they are
operational issues in the design, testing, and maintenance of modular agent
capabilities.

This paper introduces SkillOps, a practical framework for the lifecycle
management of agent skills. The term is used by analogy to operational practice
in software engineering, but the scope is narrower. SkillOps does not propose a
new model architecture or a general theory of agents. Instead, it provides a
structured way to define, test, audit, and revise skills as reusable capability
units.

The framework is derived from five open-source artifacts:

\begin{itemize}[leftmargin=*]
    \item \texttt{skill-design-guide}: guidance for defining skill structure
    and usage boundaries.
    \item \texttt{skill-security-guard}: checks for unsafe or ambiguous skill
    behavior.
    \item \texttt{persistent-memory}: a file-based memory mechanism for
    continuity across agent sessions.
    \item \texttt{agent-self-audit}: a self-auditing mechanism for reviewing
    agent outputs and operational risks.
    \item \texttt{lobster-guard}: a guardrail-oriented tool for checking
    operational assumptions and failure modes.
\end{itemize}

The central claim is operational rather than architectural: personal agent
skills become easier to inspect, test, and govern when they are treated as
lifecycle-managed artifacts with explicit structure, tests, and operational
controls. The paper contributes a framework, a taxonomy of skill components, a
methodology for artifact-based study, a descriptive benchmark layer built from
manually constructed inputs, and a local executable risk-review pilot. The
reported tables and pilot summaries should not be read as broad empirical
validation or as a general model comparison.

\section{Background and Motivation}

\subsection{From Prompts to Operational Skills}

Prompt engineering is often discussed as the design of instructions for a
language model. In personal agent systems, however, instructions are rarely
isolated. They may be tied to tools, files, context windows, memory policies,
or task-specific workflows. A skill can therefore be understood as a structured
capability unit that combines instructions with operational assumptions.

For example, a writing skill may define when it should be used, what tone it
should adopt, what sources it may rely on, and how it should handle
uncertainty. A code-review skill may specify repository access assumptions, test
execution rules, and failure-reporting requirements. A memory skill may specify
what information should be retained, where it should be stored, and how
obsolete information should be forgotten. These examples suggest that a skill is
closer to an operational module than to a single prompt.

\subsection{Operational Failure Modes}

Informal skills can fail in several recurring ways. A skill may be triggered too
broadly or too narrowly. It may inject irrelevant or stale context, causing the
model to overfit to past assumptions. It may contain constraints that are
underspecified or not testable. It may introduce safety risks when it
authorizes file edits, command execution, credential handling, or external
communication. It may also accumulate obsolete memory that conflicts with
current user intent.

These failure modes motivate a more systematic approach. If skills are treated
as artifacts, they can be linted, tested, reviewed, versioned, and retired.
This does not eliminate model uncertainty, but it can make operational
assumptions more visible and auditable. Table~\ref{tab:failure-modes}
summarizes the recurring failure categories used throughout the framework.

\begin{table}[t]
\centering
\footnotesize
\begin{tabularx}{\linewidth}{
    >{\raggedright\arraybackslash}p{0.19\linewidth}
    >{\raggedright\arraybackslash}X
    >{\raggedright\arraybackslash}p{0.31\linewidth}
}
\toprule
\textbf{Failure mode} & \textbf{Description} & \textbf{Typical mitigation} \\
\midrule
Over-triggering & The skill activates for requests that only superficially
match its scope. & Narrower trigger language, explicit negative examples,
boundary tests. \\
Under-triggering & The skill fails to activate for valid requests. & Positive
examples, paraphrase coverage, and missed-activation checks. \\
Context contamination & Irrelevant, stale, or conflicting context is injected
into execution. & Source filtering, freshness checks, conflict-resolution
rules. \\
Constraint drift & The skill's constraints are vague, contradictory, or not
enforced consistently. & Testable rules, lint checks, regression cases. \\
Unsafe tool use & The skill authorizes risky file, command, network, or
credential operations without sufficient guardrails. & Permission scoping,
security scanning, post-run audit questions. \\
Memory drift & Stored memory persists after the underlying assumption has
changed. & Expiry rules, deletion paths, explicit retirement and forgetting
policies. \\
Audit blindness & Failures are not surfaced because the skill does not require
self-review or evidence reporting. & Output contracts, self-audit prompts,
artifact logging. \\
\bottomrule
\end{tabularx}
\caption{Failure-mode taxonomy for skill-based personal agents.}
\label{tab:failure-modes}
\end{table}

\subsection{Positioning of SkillOps}

SkillOps is positioned between prompt design, tool-use agents, and software
engineering practice. It does not attempt to replace agent architectures such
as tool-use planning, reasoning-action loops, or autonomous software
engineering agents. Instead, it addresses a complementary layer: how reusable
skills should be specified and operated in a personal agent environment.

Prior work on tool-using agents studies how language models decide when and how
to use external tools \citep{schick2023toolformer,yao2023react}. Work on
agent-computer interfaces studies how environment design affects agent
performance \citep{yang2024sweagent}. Work on reusable skill libraries shows
that agent capabilities can be stored and reused over time
\citep{wang2024voyager}. SkillOps builds on this general direction but focuses
on the practical lifecycle of user-defined skills: design, triggering, context
injection, constraint enforcement, auditing, and forgetting.

\section{The SkillOps Framework}

\subsection{Definition of a Skill}

In this paper, a skill is defined as a reusable capability unit that guides an
agent's behavior for a class of tasks. A skill may include natural-language
instructions, examples, trigger conditions, context requirements, tool
permissions, safety constraints, output contracts, and tests.

SkillOps-managed skills are intended to document these assumptions in a more
structured and reviewable form. Table~\ref{tab:skill-components} summarizes the
structural components used in the framework, and Figure~\ref{fig:skill-anatomy}
provides a schematic view of how those components fit together in a skill
package.

\begin{table}[t]
\centering
\footnotesize
\begin{tabularx}{\linewidth}{
    >{\raggedright\arraybackslash}p{0.22\linewidth}
    >{\raggedright\arraybackslash}X
    >{\raggedright\arraybackslash}p{0.29\linewidth}
}
\toprule
\textbf{Component} & \textbf{Operational role} & \textbf{Example question} \\
\midrule
Name and purpose & Defines the task class and intended value of the skill. &
What problem is this skill meant to solve? \\
Trigger boundary & States when the skill should and should not activate. &
What requests are in scope, out of scope, or ambiguous? \\
Input assumptions & Records required user or environment information. & What
must be known before execution starts? \\
Context policy & Specifies what prior context, files, or memory may be
injected. & What supporting context is necessary, and how fresh must it be? \\
Execution constraints & States required, optional, and forbidden actions. &
What must the agent do, avoid, or verify? \\
Tool and file permissions & Limits resource access during execution. & Which
tools, files, networks, or services are allowed? \\
Output contract & Defines the expected form of the answer or artifact. & What
must the final output contain? \\
Safety and security checks & Surfaces risks and policy-relevant conditions. &
What unsafe actions, hidden assumptions, or trust-boundary issues must be
checked? \\
Tests and examples & Provides positive, negative, and edge cases. & How is the
skill regression-tested? \\
Lifecycle metadata & Records versioning, revision triggers, and retirement
rules. & When should the skill be updated, deprecated, or forgotten? \\
\bottomrule
\end{tabularx}
\caption{Structural components of a SkillOps-managed skill.}
\label{tab:skill-components}
\end{table}

\begin{figure}[t]
\centering
\fbox{
    \parbox{0.94\linewidth}{
        \centering
        \fbox{\parbox{0.58\linewidth}{\centering\textbf{Skill package}}}\\[0.8em]
        \setlength{\tabcolsep}{5pt}
        \renewcommand{\arraystretch}{1.05}
        \begin{tabular}{ccc}
            \diagrambox{0.25\linewidth}{\textbf{Metadata}} &
            \diagrambox{0.25\linewidth}{\textbf{Trigger}\\\textbf{contract}} &
            \diagrambox{0.25\linewidth}{\textbf{Instructions}} \\
            [0.5em]
            \diagrambox{0.25\linewidth}{\textbf{Context}\\\textbf{boundary}} &
            \diagrambox{0.25\linewidth}{\textbf{Execution}\\\textbf{constraints}} &
            \diagrambox{0.25\linewidth}{\textbf{Memory}\\\textbf{interface}} \\
            [0.5em]
            \diagrambox{0.25\linewidth}{\textbf{Tests}} &
            \diagrambox{0.25\linewidth}{\textbf{Security}\\\textbf{checks}} &
            \diagrambox{0.25\linewidth}{\textbf{Failure}\\\textbf{modes}}
        \end{tabular}
    }
}
\caption{LaTeX-native rendering of a SkillOps skill package, with the package
boundary enclosing metadata, trigger contract, instructions, context
boundary, execution constraints, memory interface, tests, security checks, and
failure modes.}
\label{fig:skill-anatomy}
\end{figure}

\subsection{Lifecycle Model}

SkillOps organizes skill management into a lifecycle:

\begin{enumerate}[leftmargin=*]
    \item \textbf{Design}: define purpose, scope, triggers, constraints, and
    output expectations.
    \item \textbf{Lint}: check for ambiguity, missing fields, conflicting
    instructions, unsafe permissions, and untestable claims.
    \item \textbf{Test}: evaluate the skill on positive, negative, boundary,
    and adversarial examples.
    \item \textbf{Deploy}: make the skill available to the agent under
    controlled activation conditions.
    \item \textbf{Observe}: collect failures, unexpected activations, user
    corrections, and context conflicts.
    \item \textbf{Audit}: review outputs and operational traces for risk,
    drift, and non-compliance.
    \item \textbf{Revise}: update the skill based on observed failures.
    \item \textbf{Forget or retire}: remove obsolete context, deprecated rules,
    or unsafe skills.
\end{enumerate}

The lifecycle emphasizes that a skill is not complete when it is written. Its
behavior depends on activation, context, model behavior, and interaction with
other skills. Figure~\ref{fig:skillops-lifecycle} summarizes the lifecycle
stages from design through forgetting or retirement.

\begin{figure}[t]
\centering
\fbox{
    \parbox{0.94\linewidth}{
        \centering
        \setlength{\tabcolsep}{4pt}
        \renewcommand{\arraystretch}{1.05}
        \begin{tabular}{c c c c c c c}
            \diagrambox{0.12\linewidth}{\textbf{Design}} &
            $\rightarrow$ &
            \diagrambox{0.12\linewidth}{\textbf{Lint}} &
            $\rightarrow$ &
            \diagrambox{0.12\linewidth}{\textbf{Test}} &
            $\rightarrow$ &
            \diagrambox{0.14\linewidth}{\textbf{Deploy}} \\
            &&&&&& $\downarrow$ \\
            \diagrambox{0.18\linewidth}{\textbf{Forget or retire}} &
            $\leftarrow$ &
            \diagrambox{0.12\linewidth}{\textbf{Revise}} &
            $\leftarrow$ &
            \diagrambox{0.12\linewidth}{\textbf{Audit}} &
            $\leftarrow$ &
            \diagrambox{0.14\linewidth}{\textbf{Observe}}
        \end{tabular}
    }
}
\caption{LaTeX-native rendering of the SkillOps lifecycle, showing the linked
sequence of design, lint, test, deploy, observe, audit, revise, and forgetting
or retirement.}
\label{fig:skillops-lifecycle}
\end{figure}

\subsection{Trigger Design}

Trigger descriptions determine when a skill is activated. A broad trigger may
widen activation coverage but also cause irrelevant activation. A narrow trigger may
reduce false positives but fail to activate in valid cases. SkillOps therefore
recommends that trigger descriptions include both positive and negative
examples.

A trigger specification should answer three questions:

\begin{itemize}[leftmargin=*]
    \item What user requests should activate the skill?
    \item What superficially similar requests should not activate it?
    \item What ambiguity should cause the agent to ask for clarification or
    fall back to a general response?
\end{itemize}

Trigger quality can be evaluated through paraphrase tests, boundary tests, and
multi-intent tests. This paper reports only the descriptive composition of a
manually constructed trigger benchmark; Section~\ref{sec:evaluation}
summarizes the reported counts and their limitations.

\subsection{Context Injection}

Context injection is necessary when a skill depends on project conventions,
user preferences, prior decisions, or repository structure. However, excessive
context can degrade behavior by introducing stale assumptions or irrelevant
constraints. SkillOps treats context injection as a declared policy rather than
opportunistic accumulation of memory.

A context injection policy should specify:

\begin{itemize}[leftmargin=*]
    \item the minimum context needed for the skill to operate;
    \item the source of the context, such as files, memory entries, or tool
    outputs;
    \item the freshness requirements for the context;
    \item the conditions under which context should be omitted;
    \item the method for resolving conflicts between current user instructions
    and stored context.
\end{itemize}

This is especially important for personal agents, where long-lived memory can
improve continuity but also preserve outdated assumptions.

\subsection{Execution Constraints}

Execution constraints define what the agent is allowed to do during skill use.
For low-risk writing skills, constraints may concern tone, citation style, or
formatting. For higher-risk skills, constraints may concern file modification,
command execution, external communication, credential handling, or irreversible
actions.

SkillOps distinguishes between three types of constraints:

\begin{itemize}[leftmargin=*]
    \item \textbf{Behavioral constraints}: rules about tone, reasoning style,
    response structure, or uncertainty reporting.
    \item \textbf{Operational constraints}: rules about tools, files, commands,
    commits, network access, and external services.
    \item \textbf{Safety constraints}: rules about privacy, credentials,
    harmful actions, or irreversible operations.
\end{itemize}

A constraint is more useful when it is testable. For example, ``do not claim
success without evidence'' is more testable than ``be careful.'' SkillOps
therefore favors constraints that can be checked through examples, lint rules,
or audit questions.

\subsection{Forgetting and Retirement}

Forgetting is a necessary part of skill operation. A personal agent may retain
preferences, project state, or procedural rules that later become obsolete.
Without a forgetting mechanism, the agent may apply outdated instructions to
new tasks.

SkillOps treats forgetting at two levels. First, individual memory entries may
need to be removed or revised. Second, entire skills may need to be deprecated
when their scope changes or when they become unsafe. Retirement rules should
specify how to identify obsolete skills, how to preserve historical context if
needed, and how to prevent deprecated instructions from continuing to influence
behavior.

Table~\ref{tab:informal-vs-skillops} contrasts informal prompt-based skills with
skills managed under the proposed framework.

\begin{table}[t]
\centering
\footnotesize
\begin{tabularx}{\linewidth}{
    >{\raggedright\arraybackslash}p{0.23\linewidth}
    >{\raggedright\arraybackslash}p{0.34\linewidth}
    >{\raggedright\arraybackslash}X
}
\toprule
\textbf{Dimension} & \textbf{Informal skill} & \textbf{SkillOps-managed skill} \\
\midrule
Scope definition & Implicit task description embedded in prompt text. & Named
purpose, explicit scope, and documented non-scope. \\
Activation & Heuristic or ad hoc triggering. & Trigger contract with positive,
negative, and boundary examples. \\
Context handling & Opportunistic context accumulation. & Declared context
sources, freshness rules, and conflict handling. \\
Constraints & Broad instructions such as ``be careful.'' & Testable behavioral,
operational, and safety constraints. \\
Testing & Rare or manual spot checks. & Regression cases, adversarial examples,
and lint coverage. \\
Security review & Often absent or implicit. & Permission scoping, unsafe-action
checks, and trust-boundary review. \\
Memory handling & Persistent context without explicit expiry rules. & Defined
write permissions, forgetting policy, and retirement criteria. \\
Operational reporting & Minimal trace of what changed and why. & Output
contract, audit fields, versioning, and revision history. \\
\bottomrule
\end{tabularx}
\caption{Comparison between informal prompt-based skills and SkillOps-managed
skills.}
\label{tab:informal-vs-skillops}
\end{table}

\subsection{Problem Formulation and Design Gap}

The core problem can be stated as follows. Given a set of reusable skills
\(\mathcal{S}\), a user request \(x\), optional context \(c\), and an execution
environment \(e\), an agent must decide whether to activate a skill, which
context to expose, which operations to allow, and how to report uncertainty or
failure. Informal prompt-based skills usually compress these decisions into a
single natural-language instruction. SkillOps separates them into reviewable
interfaces: activation, context, constraints, memory, and audit.

This separation matters because the failure modes are coupled. An over-broad
trigger can expose irrelevant memory; stale memory can defeat otherwise clear
constraints; weak output contracts can hide unsafe tool use; and missing
retirement rules can keep obsolete behavior alive after a project changes.
Table~\ref{tab:design-gap} summarizes the design gap addressed by the
framework.

\begin{table}[t]
\centering
\footnotesize
\begin{tabularx}{\linewidth}{
    >{\raggedright\arraybackslash}p{0.17\linewidth}
    >{\raggedright\arraybackslash}p{0.27\linewidth}
    >{\raggedright\arraybackslash}X
}
\toprule
\textbf{Axis} & \textbf{Informal handling} & \textbf{SkillOps interface} \\
\midrule
Activation & Broad task description or keyword matching. & Trigger contract
with positive, negative, and ambiguous cases. \\
Context & Whatever prior state is convenient to include. & Declared sources,
freshness rules, and conflict handling. \\
Constraints & General cautionary language. & Testable behavioral,
operational, and safety constraints. \\
Memory & Persistent context without retirement semantics. & Explicit write,
read, forgetting, and deprecation policy. \\
Audit & Success or failure reported after the fact. & Output contract,
evidence fields, and review checks tied to known failure modes. \\
\bottomrule
\end{tabularx}
\caption{SkillOps turns coupled skill-design risks into explicit interfaces
that can be reviewed and regression-tested.}
\label{tab:design-gap}
\end{table}

\section{Research Questions}

This paper is organized around three research questions.

\paragraph{RQ1: What structural components should a skill include as an agent capability unit?}

This question examines whether a practical skill template can make agent
capabilities easier to understand, test, and maintain. The expected output is a
component taxonomy covering triggers, context, constraints, tool access, output
contracts, tests, and lifecycle metadata.

\paragraph{RQ2: How do trigger descriptions, context injection, execution
constraints, and forgetting expose operational risk?}

This question examines how design choices influence activation behavior,
consistency across paraphrases, constraint following, and resistance to stale
context. The expected output is a benchmark suite and a provider-neutral
experimental protocol with controlled variations of skill definitions.

\paragraph{RQ3: How can linting, security scanning, and self-auditing support
operational risk review through reviewable checks?}

This question examines how reviewable checks can surface common skill risks and
make them easier to inspect. The expected output is a manually constructed set
of lint, security, and self-audit cases together with repository-level checks
and review criteria for operational risk review. It does not by itself prove
that these mechanisms reduce deployment risk or improve model behavior.

Table~\ref{tab:rq-mapping} maps each research question to the corresponding
framework component and evaluation method.

\begin{table}[t]
\centering
\footnotesize
\begin{tabularx}{\linewidth}{
    >{\raggedright\arraybackslash}p{0.18\linewidth}
    >{\raggedright\arraybackslash}p{0.29\linewidth}
    >{\raggedright\arraybackslash}X
}
\toprule
\textbf{Research question} & \textbf{Framework component} &
\textbf{Evaluation method} \\
\midrule
RQ1 & Skill structure and documentation fields & Artifact review, template
comparison, and completeness analysis across the open-source repositories. \\
RQ2 & Trigger design, context policy, execution constraints, and forgetting &
Manually constructed benchmark cases and model-backed protocols covering
activation boundaries, stale-context injections, and deprecated-memory
scenarios. \\
RQ3 & Linting, security scanning, and self-audit & Seeded failure cases,
reviewable repository checks, and manual review criteria for operational risk
inspection. \\
\bottomrule
\end{tabularx}
\caption{Mapping from research questions to framework components and evaluation
methods.}
\label{tab:rq-mapping}
\end{table}

\section{Artifacts and Methodology}

\subsection{Artifact Base}

The framework is derived from a set of open-source artifacts.
Table~\ref{tab:artifacts} summarizes their intended roles.

\begin{table}[t]
\centering
\footnotesize
\begin{tabularx}{\linewidth}{
    >{\raggedright\arraybackslash}p{0.29\linewidth}
    >{\raggedright\arraybackslash}X
}
\toprule
\textbf{Artifact} & \textbf{Role in SkillOps} \\
\midrule
\texttt{skill-design-guide} & Provides design guidance for structuring skills,
defining scope, and writing trigger descriptions. \\
\texttt{skill-security-guard} & Provides checks for unsafe, ambiguous, or
overly permissive skill behavior. \\
\texttt{persistent-memory} & Provides a file-based approach to retaining and
distilling long-term agent context. \\
\texttt{agent-self-audit} & Provides a mechanism for reviewing agent outputs,
assumptions, and possible failure modes. \\
\texttt{lobster-guard} & Provides operational guardrails for checking
assumptions, reporting partial failure, and avoiding unsupported claims. \\
\bottomrule
\end{tabularx}
\caption{Open-source artifacts used as the practical basis for the SkillOps
framework.}
\label{tab:artifacts}
\end{table}

The repository location for the paper and its supporting materials is provided
in Section~\ref{sec:artifact-availability}.

\subsection{Methodological Approach}

This paper uses an artifact-based methodology. Rather than beginning with a
large-scale user study, it derives a framework from the design and use of
practical open-source tools. The methodology has four steps:

\begin{enumerate}[leftmargin=*]
    \item \textbf{Artifact review}: identify recurring structures and
    constraints across the repositories.
    \item \textbf{Component abstraction}: generalize these structures into a
    skill component taxonomy.
    \item \textbf{Failure-mode analysis}: identify common ways in which skills
    can misfire, become stale, or create operational risk.
    \item \textbf{Design-gap mapping}: map each failure mode to the skill
    interface that should make it inspectable.
    \item \textbf{Benchmark design}: construct benchmark cases intended for
    later evaluation of trigger behavior, context injection, constraint
    following, forgetting, linting, security scanning, and self-audit
    behavior.
\end{enumerate}

This approach is exploratory. It is appropriate for deriving a
framework, but it cannot by itself establish broad generality. The limitations
are discussed in Section~\ref{sec:limitations}.

\subsection{Benchmark Construction and Descriptive Outputs}

The benchmark includes manually constructed cases in the following
categories:

\begin{itemize}[leftmargin=*]
    \item \textbf{Activation cases}: requests that should activate a skill.
    \item \textbf{Non-activation cases}: requests that appear related but should
    not activate the skill.
    \item \textbf{Boundary cases}: ambiguous requests requiring clarification or
    fallback.
    \item \textbf{Context cases}: tasks requiring relevant context and tasks
    where injected context is harmful.
    \item \textbf{Constraint cases}: tasks testing whether operational and
    safety constraints are followed.
    \item \textbf{Forgetting cases}: tasks where obsolete memory or deprecated
    instructions should not be used.
    \item \textbf{Security cases}: tasks containing unsafe permissions,
    credential exposure, prompt injection, or irreversible actions.
    \item \textbf{Audit cases}: completed outputs requiring self-review for
    unsupported claims, missing evidence, or hidden assumptions.
\end{itemize}

The repository includes manually constructed benchmark inputs and
generated descriptive outputs. \path{benchmark/skill_samples.csv} profiles
five artifacts, \path{benchmark/trigger_cases.csv} contains 36 trigger
cases, and \path{benchmark/risk_cases.csv} contains 24 operational-risk
cases. Running \path{scripts/run_all.py} regenerates summary tables under
\path{results/tables/} and SVG figures under \path{figures/}. Because the
benchmark is manually constructed and exploratory, these outputs should be read
as descriptive summaries rather than as statistical evidence or broad empirical
validation. The repository also includes executable repository-level
consistency checks that verify benchmark schema, result reproducibility,
paper-claim consistency, public-presentation constraints, SVG validity, and
evidence-matrix integrity. In the current repository state, this executable
pass regenerates descriptive tables and schematic figures only; it does not
execute the five source repositories against the benchmark cases or score
model, scanner, or self-audit behavior.

\subsection{Artifact Availability}
\label{sec:artifact-availability}

The public repository for this paper is
\url{https://github.com/rrrrrredy/skillops-paper}.

Supporting materials are organized as follows:
\begin{itemize}[leftmargin=*]
    \item Artifact inventory: \path{artifacts/artifact_inventory.md}
    \item Benchmark input: \path{benchmark/skill_samples.csv}
    \item Benchmark input: \path{benchmark/trigger_cases.csv}
    \item Benchmark input: \path{benchmark/risk_cases.csv}
    \item Reproducible scripts: \path{scripts/}
    \item Preserved pilot reports and metrics: \path{results/experiments/}
    \item Generated result tables: \path{results/tables/}
    \item Generated figure artwork: \path{figures/}
\end{itemize}

These materials are intended to make the exploratory benchmark inspectable,
but the benchmark remains manually constructed and should not be interpreted as
broad empirical validation.

\section{Evaluation}
\label{sec:evaluation}

This revision reports three evidence layers currently versioned in the
repository. The first is a descriptive layer derived from manually constructed
benchmark artifacts. The second is a deterministic repository-check layer. The
third is a local security-guard pilot over seeded risk cases and benign
controls. A model-backed protocol is specified in the experiment harness, but no
model-backed live metrics are reported in this version. No statistical
significance is claimed, no broad empirical validation is claimed, and the
results should be read as evidence of inspectability and executable review
rather than proof that SkillOps is generally effective.

\subsection{Evaluation Scope}

The evaluation combines the following versioned inputs:
\begin{itemize}[leftmargin=*]
    \item Skill samples: \path{benchmark/skill_samples.csv}
    \item Trigger cases: \path{benchmark/trigger_cases.csv}
    \item Risk cases: \path{benchmark/risk_cases.csv}
    \item Benign controls: \path{experiments/security_benign_cases.csv}
    \item Cross-checking notes: \path{artifacts/artifact_inventory.md}
    \item Local pilot metrics: \path{results/experiments/security_guard_metrics.csv}
\end{itemize}

These inputs cover five artifact profiles, 36 trigger-routing cases, 24
operational-risk cases, and 24 benign security controls. All trigger and risk
cases use the \texttt{manual-} prefix, making the manual construction of the
benchmark explicit. Figure~\ref{fig:evaluation-pipeline} captures the
evaluation flow from artifact inventory to cautious interpretation.

\begin{figure}[t]
\centering
\fbox{
    \parbox{0.94\linewidth}{
        \centering
        \setlength{\tabcolsep}{4pt}
        \renewcommand{\arraystretch}{1.05}
        \begin{tabular}{c c c c c}
            \diagrambox{0.21\linewidth}{\textbf{Artifact}\\\textbf{inventory}} &
            $\rightarrow$ &
            \diagrambox{0.21\linewidth}{\textbf{Manual benchmark}\\\textbf{cases}} &
            $\rightarrow$ &
            \diagrambox{0.21\linewidth}{\textbf{Executable}\\\textbf{checks}} \\
            &&&& $\downarrow$ \\
            \diagrambox{0.21\linewidth}{\textbf{Cautious}\\\textbf{interpretation}} &
            $\leftarrow$ &
            \diagrambox{0.21\linewidth}{\textbf{Result}\\\textbf{tables}} &
            $\leftarrow$ &
            \diagrambox{0.21\linewidth}{\textbf{Local}\\\textbf{pilot}}
        \end{tabular}
    }
}
\caption{Evaluation pipeline from artifact inventory and manual benchmark cases
through executable checks, result tables, a local pilot, and cautious
interpretation.}
\label{fig:evaluation-pipeline}
\end{figure}

\subsection{Descriptive Benchmark Summaries}

The descriptive benchmark layer documents benchmark composition and artifact
coding only. Running \path{scripts/run_all.py} regenerates summary tables
under \path{results/tables/} and SVG figures under \path{figures/}, but
it does not execute the five source repositories or score routing, linting,
scanning, or self-audit behavior.

\paragraph{Artifact coverage.}

Table~\ref{tab:artifact-coverage-summary} summarizes the descriptive coding of
nine SkillOps-relevant components across the five inspected artifacts.
Metadata, trigger contracts, instructions, context boundaries, execution
constraints, security checks, and failure modes are documented in all five
artifacts. Memory interfaces are documented in one artifact and limited in
four, while tests are documented in three artifacts and limited in two.

\begin{table}[t]
\centering
\footnotesize
\begin{tabular}{lrrr}
\toprule
\textbf{Component} & \textbf{Documented} & \textbf{Limited} & \textbf{Absent} \\
\midrule
Metadata & 5 & 0 & 0 \\
Trigger contract & 5 & 0 & 0 \\
Instructions & 5 & 0 & 0 \\
Context boundary & 5 & 0 & 0 \\
Execution constraints & 5 & 0 & 0 \\
Memory interface & 1 & 4 & 0 \\
Tests & 3 & 2 & 0 \\
Security checks & 5 & 0 & 0 \\
Failure modes & 5 & 0 & 0 \\
\bottomrule
\end{tabular}
\caption{Descriptive artifact-coverage counts generated from the manually
constructed artifact profiles.}
\label{tab:artifact-coverage-summary}
\end{table}

\paragraph{Trigger cases.}

The trigger benchmark contains 36 manually written cases: 15 positive
activation cases, 12 negative cases, and 9 ambiguous cases. The cases are
distributed across the five inspected skills plus seven no-route cases that are
intended not to activate any inspected skill. Table~\ref{tab:trigger-summary}
reports these counts directly from the generated summary files.

\begin{table}[t]
\centering
\footnotesize
\begin{tabularx}{\linewidth}{
    >{\raggedright\arraybackslash}p{0.24\linewidth}
    >{\raggedright\arraybackslash}X
    >{\raggedleft\arraybackslash}p{0.12\linewidth}
}
\toprule
\textbf{Dimension} & \textbf{Category} & \textbf{Count} \\
\midrule
Expected label & \texttt{should\_trigger} & 15 \\
Expected label & \texttt{should\_not\_trigger} & 12 \\
Expected label & \texttt{ambiguous} & 9 \\
Relevant skill & \texttt{skill-design-guide} & 7 \\
Relevant skill & \texttt{skill-security-guard} & 6 \\
Relevant skill & \texttt{persistent-memory} & 7 \\
Relevant skill & \texttt{agent-self-audit} & 4 \\
Relevant skill & \texttt{lobster-guard} & 5 \\
Relevant skill & \texttt{none} & 7 \\
\bottomrule
\end{tabularx}
\caption{Descriptive trigger-case counts generated from the manually
constructed trigger benchmark.}
\label{tab:trigger-summary}
\end{table}

\paragraph{Risk cases.}

The risk benchmark contains 24 manually written cases. The case inventory is
evenly split across eight risk types, with three cases each for prompt
injection, over-broad triggers, unsafe file access, missing constraints, stale
memory, missing tests, identity confusion, and privacy leakage.
Table~\ref{tab:risk-summary} also reports the artifact emphasis of the case
set, with six cases tied most directly to \texttt{persistent-memory}, five
each to \texttt{skill-security-guard} and \texttt{lobster-guard}, and four
each to \texttt{skill-design-guide} and \texttt{agent-self-audit}.

\begin{table}[t]
\centering
\footnotesize
\begin{tabularx}{\linewidth}{
    >{\raggedright\arraybackslash}p{0.24\linewidth}
    >{\raggedright\arraybackslash}X
    >{\raggedleft\arraybackslash}p{0.12\linewidth}
}
\toprule
\textbf{Dimension} & \textbf{Category} & \textbf{Count} \\
\midrule
Risk type & \texttt{prompt\_injection} & 3 \\
Risk type & \texttt{over\_broad\_trigger} & 3 \\
Risk type & \texttt{unsafe\_file\_access} & 3 \\
Risk type & \texttt{missing\_constraints} & 3 \\
Risk type & \texttt{stale\_memory} & 3 \\
Risk type & \texttt{missing\_tests} & 3 \\
Risk type & \texttt{identity\_confusion} & 3 \\
Risk type & \texttt{privacy\_leakage} & 3 \\
Relevant artifact & \texttt{skill-security-guard} & 5 \\
Relevant artifact & \texttt{lobster-guard} & 5 \\
Relevant artifact & \texttt{skill-design-guide} & 4 \\
Relevant artifact & \texttt{persistent-memory} & 6 \\
Relevant artifact & \texttt{agent-self-audit} & 4 \\
\bottomrule
\end{tabularx}
\caption{Descriptive risk-case counts generated from the manually constructed
operational-risk benchmark.}
\label{tab:risk-summary}
\end{table}

\subsection{Repository-Level Consistency Checks}

The repository includes \path{scripts/run_tests.py}, a deterministic
repository-level test suite for consistency and traceability checks. In the
current repository state, the suite covers benchmark schema, result
reproducibility, paper-claim consistency, public-presentation constraints, SVG
validity, and evidence-matrix integrity, and it passes the full repository
test suite.

These checks validate repository consistency and traceability only. They do not
measure model behavior, scanner accuracy, user-study outcomes, production
validation, or statistical significance.

\subsection{Local Security-Guard Pilot}

The local pilot executes \path{scripts/run_security_guard_experiment.py} with
the \path{local-rules} guard over the 24 risk cases and 24 benign controls. It
is intentionally simple: the guard is a deterministic pattern-based checker, so
the pilot evaluates whether the seeded risk corpus and benign controls are
wired into an executable review loop, not whether a learned model can identify
new risks.

\begin{table}[t]
\centering
\footnotesize
\begin{tabular}{lrr}
\toprule
\textbf{Metric} & \textbf{Value} & \textbf{Count} \\
\midrule
Risk detection rate & 1.000000 & 24/24 \\
Benign false-positive rate & 0.041667 & 1/24 \\
Benign specificity & 0.958333 & 23/24 \\
\bottomrule
\end{tabular}
\caption{The local security-guard pilot caught all seeded risk cases and
flagged one benign control, showing executable risk-review plumbing rather than
broad security validation.}
\label{tab:local-security-pilot}
\end{table}

All eight risk categories had three of three seeded cases detected in this
local run. The single benign false positive is useful rather than incidental:
it indicates that guard checks need benign controls whenever the risk detector
uses broad textual patterns.

\subsection{Model-Backed Experiment Protocol}

The repository also specifies model-backed live-run protocols for trigger
routing, constraint handling, model-backed security review, memory drift, and
skill-component ablations. The protocol fixes input files, prompt templates,
JSON schemas, metric calculations, and dry-run validation before any live call.
In this repository state, readiness and dry-run checks pass, but no
model-backed live metrics are reported. A future run should record the selected
model, exact date, raw outputs, parse failures, execution failures, and
confidence boundaries before any result is cited in the paper.

\subsection{Interpretation}

Taken together, the descriptive tables document benchmark composition, the
repository checks validate internal consistency, and the local security pilot
shows that at least one review loop can execute end to end over risk and benign
cases. The strongest supported claim is therefore not that SkillOps improves
agent performance. The supported claim is narrower and, for this paper, more
important: explicit skill interfaces make operational assumptions, failure
modes, and review checks concrete enough to inspect, test, and revise.

These results do not establish statistical significance, broad empirical
validation, or general effectiveness across external skill sets, user
populations, or deployment environments.

\section{Discussion}

SkillOps frames personal agent skills as operational artifacts rather than
isolated prompts. This framing has several implications.

The main practical insight is not that SkillOps makes skill behavior
deterministic. Rather, it moves trigger cases, context policies, constraints,
and evidence logs into reviewable artifacts that support measurable checks and
reviewable operational analysis.

First, it makes skill design easier to review. When trigger conditions, context
policies, and constraints are written clearly, they can be reviewed by humans
and checked by tools. This is useful even when the underlying model remains
uncertain.

Second, it encourages test-driven skill development. A skill should include
examples of intended activation, unintended activation, edge cases, and unsafe
requests. These examples do not guarantee correct behavior, but they create a
practical basis for regression testing.

Third, it treats memory as both useful and risky. Persistent memory can improve
continuity across sessions, but it can also preserve outdated preferences,
project state, or unsafe assumptions. SkillOps therefore treats forgetting as a
first-class operation rather than an afterthought.

Fourth, it suggests that safety checks should be integrated into normal skill
operation. Security scanning and self-auditing should not be reserved only for
high-risk deployments. Even personal agents can perform file edits, repository
operations, or external communications that require explicit safeguards.

Finally, SkillOps may help bridge individual practice and more formal agent
engineering. Many agent failures are not caused solely by model limitations.
They also arise from unclear instructions, uncontrolled context, weak
operational boundaries, and missing review procedures. A lightweight
operational framework can make these issues easier to identify and improve.

\section{Limitations}
\label{sec:limitations}

This paper has several limitations.

First, the artifact base is authored and interpreted by one author. This paper
therefore reflects a single-author artifact base
rather than a multi-lab or multi-organization corpus.

Second, the artifact evidence comes from five public repositories. This is
enough to motivate a framework, but it is not enough to establish generality
across the broader agent ecosystem.

Third, the benchmark cases are manually constructed and exploratory. The counts
reported in Section~\ref{sec:evaluation} describe benchmark composition rather
than observed production traffic or a statistically powered experiment.

Fourth, the main executable pilot in this version is local and rule-based. It
checks the risk-review plumbing, but it is not evidence that a model-backed
agent will make the same decisions.

Fifth, the benchmark cases include manually constructed benign security
controls. The reported detection and false-positive rates are therefore
benchmark-specific rather than broad external validation.

Sixth, model-backed trigger, constraint, memory, and ablation experiments remain
specified but not reported as live metrics in this version. Any future report
should pre-register the selected model, prompt hashes, case set, repeat count,
and metric calculation before interpreting the results.

Seventh, no statistical significance or formal hypothesis testing is reported,
and external validation remains limited.

Eighth, the work does not include an external user study. It does not measure
how other users understand the skill format, apply the workflow, or experience
the proposed guardrails in practice.

Ninth, external skill sets, other agent environments, and independent
annotations are needed before stronger comparative claims could be made.

Tenth, several artifacts encode local desktop-agent assumptions about directory
layout, skill packaging, cron usage, and local agent control. These assumptions
may not transfer directly to hosted agents, enterprise workflows, or other
agent environments.

Eleventh, SkillOps focuses on the operational layer of skill design. It does not
propose a new model architecture, training method, planner, or tool-use
algorithm.

Twelfth, some evaluation criteria still require human judgment. For example,
audit usefulness and context relevance may be difficult to score objectively.
Later benchmark revisions should therefore include clear annotation guidelines.

Thirteenth, personal agent environments differ substantially in their tool access,
memory mechanisms, security policies, and user expectations. A skill format
that works in one environment may require adaptation in another.

\section{Related Work}

\paragraph{Agent skills and skill-library maintenance.}

The recent Agent Skills ecosystem makes the skill abstraction a first-class
interface rather than a private prompt convention. Anthropic's Agent Skills
package procedural knowledge in file-and-folder bundles with progressive
context disclosure and explicit security cautions
\citep{anthropic2025agentskills}. SkillsBench evaluates curated and
self-generated skills across many tasks and finds that skill usefulness is
large but uneven across domains, which reinforces the need to test skills as
artifacts rather than assume that more procedural context is always beneficial
\citep{li2026skillsbench}. A recent paper using the same SkillOps name studies
library-level maintenance through typed skill contracts and a hierarchical skill
ecosystem graph \citep{pu2026skillops}. This paper is narrower in one respect
and broader in another: it does not propose a library-maintenance algorithm or
report environment task success, but it treats personal-agent skills as
operational artifacts whose triggers, context, constraints, memory, tests, and
audit rules must be inspectable even before a large skill library exists.

\paragraph{Tool-using and modular agents.}

SkillOps is discussed here mainly as a framework for agent skill artifacts
within work on modular agents and tool-using language model systems, with
earlier software-modularity references serving only as background analogies.
Classical work on software decomposition and reusable structure emphasized
information hiding, design patterns, and component-based modularity
\citep{parnas1972criteria,gamma1994design,szyperski2002component}, while
hierarchical reinforcement learning formalized reusable temporally extended
skills \citep{sutton1999options,dietterich2000maxq}. In the LLM era, MRKL,
Toolformer, and ReAct study how models call tools or coordinate reasoning with
external actions \citep{karpas2022mrkl,schick2023toolformer,yao2023react}.
Benchmark-oriented tool-use work such as API-Bank and ToolLLM evaluates API
invocation over larger tool sets \citep{li2023apibank,qin2024toolllm}, while
SayCan and Gorilla connect language models to robotic affordances or large API
spaces \citep{ichter2023saycan,patil2024gorilla}. LATM and CRAFT study
tool-making or toolset customization \citep{cai2024latm,yuan2024craft}, and
AutoGen and MetaGPT organize multi-agent software workflows
\citep{wu2024autogen,hong2024metagpt}. SkillOps differs by treating a skill as
a governed artifact with trigger contracts, context policy, constraints, tests,
and lifecycle metadata rather than only as an action primitive, generated tool,
or coordination role.

\paragraph{Skill libraries and memory.}

Reusable skill libraries appear in embodied and long-horizon agents such as
Voyager, where executable skills are stored and reused over time
\citep{wang2024voyager}. Generative Agents, CoALA, MemoryBank, and MemGPT show
that long-running agents need explicit mechanisms for memory storage,
retrieval, reflection, and control
\citep{park2023generativeagents,sumers2024coala,zhong2024memorybank,packer2023memgpt}.
Work on reflective improvement, including Reflexion and Self-Refine, further
suggests that agent behavior can be improved by structured feedback loops
\citep{shinn2023reflexion,madaan2023selfrefine}. SkillOps extends this line by
specifying which skills may write or consume memory, how staleness is handled,
and how retirement or forgetting is enforced.

\paragraph{Agent security and prompt injection.}

Security work on LLM-integrated applications highlights why skill-level context
policy matters. Indirect prompt injection shows that untrusted content can
override intended instructions when data and control are mixed
\citep{greshake2023notwhat}. Red-teaming, formal prompt-injection benchmarking,
instruction-priority training, structured-query defenses, and environments such
as AgentDojo all study ways to characterize or mitigate these risks
\citep{ganguli2022redteaming,liu2024formalizingprompt,wallace2024instructionhierarchy,chen2025struq,debenedetti2024agentdojo}.
SkillOps does not claim to solve prompt injection; its narrower goal is to
localize trust boundaries, permission scopes, memory-write rules, and audit
fields within each skill.

\paragraph{Agent evaluation and software engineering.}

Recent agent evaluation work increasingly uses interactive or executable
environments rather than static question answering. SWE-bench and SWE-agent
study real software tasks and the importance of agent-computer interfaces
\citep{jimenez2024swebench,yang2024sweagent}, while AgentBench and
\(\tau\)-bench broaden evaluation to multiple environments and tool-user
interaction settings \citep{liu2024agentbench,yao2025taubench}.
SkillOps adopts the operational spirit of this literature but focuses on
skill-level properties such as false-trigger rate, stale-context errors,
constraint violations, and recovery after partial failure.

\paragraph{Documentation, model cards, and operational reporting.}

SkillOps is also informed by work on documentation and operational discipline
for AI systems. Large language models and instruction tuning make
natural-language control surfaces central to system behavior
\citep{brown2020language,ouyang2022training}, while human-AI interaction
guidelines emphasize communicating system capabilities, uncertainty, and user
expectations \citep{amershi2019guidelines}. Model Cards, Datasheets, Hidden
Technical Debt, and the ML Test Score argue that AI systems need structured
reporting, monitoring, and readiness checks
\citep{mitchell2019modelcards,gebru2021datasheets,sculley2015hidden,breck2017mltestscore}.
SkillOps transfers this documentation mindset to modular agent skills by making
operational assumptions explicit and reviewable.

\section{Conclusion}

This paper introduced SkillOps, a practical framework for designing, testing,
and operating modular skills in personal AI agents. The framework is derived
from open-source artifacts and emphasizes explicit skill structure, trigger
design, context injection, execution constraints, security checks,
self-auditing, and forgetting. The paper formulates three research questions
and reports a descriptive benchmark layer, repository-level consistency checks,
and a local security-guard pilot over risk and benign-control cases.

The contribution is bounded but operationally substantive. SkillOps is not
presented as a new agent architecture or as a broadly validated empirical
result. Instead, it identifies an under-specified layer in personal-agent
systems: the lifecycle discipline required to make reusable skills inspectable,
testable, and maintainable. The controlled pilot shows that structured skill
artifacts can support executable checks and reviewable operational analysis in
this setting, while broader empirical follow-up would require external skill
sets, model-backed runs with preserved raw logs, repeated execution studies,
and refinement based on external use.

\bibliographystyle{plainnat}
\bibliography{references}

\end{document}
